\begin{textsample}{POS Dim 4 – rr – Score 12.00 – rr/rr\_em\_15.txt}  \label{ex:f4_pos_066}
2.2.
O Estado de Roraima, localizado no extremo \textbf{norte} do Brasil, concentra uma grande diversidade étnica de povos tradicionais da América, historicamente conhecidos como \textbf{índios}.
Desde o Período Colonial até aos dias atuais, o contato dos povos \textbf{indígenas} com a sociedade não \textbf{indígena} nacional, tem produzido avanços e retrocessos no sentido de conservação e respeito às diferenças culturais em seus mais diversos aspectos.
Foi a partir de 1988, que a Constituição Federal ( BRASIL, 1988 ), dispensou a orientação legal para a prática da educação escolar \textbf{indígena} no país.
Nesse sentido, o art. 231 orienta para a proteção da cultura, língua e tradições \textbf{indígenas}.
Segundo dados da Divisão de Educação Escolar \textbf{Indígena} – DIEI/DGE/SEED/RR, na rede estadual de ensino há um total de 264 escolas \textbf{indígenas}, distribuídas em 12 regiões, das quais, 71 ofertam Ensino Médio, conforme apresentado no quadro abaixo : | Região | Total de Escolas \textbf{Indígenas} da Rede Estadual | Nº de Escolas \textbf{Indígenas} que ofertam Ensino Médio | |---|---:|---:| | Amajari | 16 | 06 | | Baixo Cotingo | 27 | 08 | | Murupú | 03 | 02 | | Raposa | 26 | 06 | | São Marcos | 31 | 13 | | Serra da Lua | 21 | 09 | | Serras | 65 | 16 | | Surumu | 23 | 04 | | Tabaio | 10 | 05 | | Wai Wai | 08 | 02 | | Yanomami e Yekuna | 34 | - | | Total | 264 | 71 | Fonte : Divisão de Educação Escolar \textbf{Indígena} – DIEI/DGE/SEED/RR.
Assim como os \textbf{indígenas}, os ribeirinhos também são considerados povos tradicionais que habitam as margens dos \textbf{rios} na imensa \textbf{floresta} \textbf{amazônica}, praticam a agricultura, a pesca artesanal e o \textbf{artesanato} como meio principal de subsistência.
Esses povos possuem uma cultura bem diferenciada de outros grupos sociais e suas habitações são típicas, adaptadas para as estações chuvosas em que a vazão dos \textbf{rios} aumenta, colocando em risco as construções das casas em forma de palafitas.
O povo ribeirinho, tal qual o povo \textbf{indígena}, também tem os seus direitos legais assegurados na Constituição Federal ( BRASIL, 1988 ) e em outras normativas ( BRASIL, 1995 ; BRASIL, 1996 ), contando com um modelo de educação específica para o desenvolvimento dos processos pedagógicos da cultura ribeirinha.
As comunidades ribeirinhas no Estado de Roraima, localizam -se na região do Baixo Rio Branco, \textbf{sul} do Estado, fazendo fronteira com o Estado do Amazonas.
Dentre as muitas vilas, destacam -se : Santa Maria do Boiaçú, Santa Maria Velha, Terra Preta, Sacaí, Lago Grande, Caicumbí, \textbf{Floresta}, Itaquera, Samaúma, Panacarica e Remanso.
Para essas \textbf{populações} presentes no estado de Roraima, o Novo Ensino Médio deve ser implementado considerando as legislações específicas que asseguram os direitos desses sujeitos.
Já a \textbf{população} rural, que abrange um público diverso, deve ser atendida pela modalidade de educação do campo, tal qual considerado pelo Parecer CNE/CEB nº 36 ( BRASIL, 2001, p.01 ) : > A educação do campo, tratada como educação rural na legislação brasileira, tem um significado que incorpora os espaços da \textbf{floresta}, da pecuária, das minas e da agricultura, mas os ultrapassa ao acolher em si, os espaços pesqueiros, caiçara, ribeirinho e extrativistas.
Nesse sentido, os sujeitos do campo ocupam mais que um espaço não-urbano, o qual deve ser entendido pela escola como um lugar de possibilidades para dinamizar a ligação dos seres humanos com a própria produção das condições de existência social e com as realizações da sociedade humana.
Portanto, a escola em seu Projeto Pedagógico deve incorporar a concepção da educação integral considerando a diversidade e especificidade dos povos tradicionais.

% matched lemmas: amazônico, artesanato, floresta, indígena, norte, população, rio, sul, índio
\end{textsample}
